% Réécriture en flux de l'export Excel général de mod_stage (export.php) : méthode, équivalence
% du classeur produit, durée et mémoire avant et après.
% Compilation : latexmk -pdf export-en-flux.tex
\documentclass[11pt,a4paper]{article}
\usepackage[french]{babel}
\usepackage[T1]{fontenc}
\usepackage{lmodern,microtype,booktabs,tabularx,xcolor,hyperref,enumitem}
\usepackage[margin=2.2cm]{geometry}
\usepackage{url}
\expandafter\def\expandafter\UrlBreaks\expandafter{\UrlBreaks\do\a\do\b\do\c\do\d\do\e\do\f\do\g\do\h\do\i\do\j\do\k\do\l\do\m\do\n\do\o\do\p\do\q\do\r\do\s\do\t\do\u\do\v\do\w\do\x\do\y\do\z}
\DeclareUrlCommand\code{\urlstyle{tt}}
\Urlmuskip=0mu plus 2mu
\setlist[itemize]{itemsep=2pt,topsep=3pt}
\title{Export Excel général de \texttt{mod\_stage} : écriture en flux}
\author{Sébastien Lefebvre}
\date{11 novembre 2026}

\begin{document}
\maketitle

\section{Objet}

L'audit des types de question (\code{audit-types-de-question.tex}) signalait la lenteur de l'export
Excel général, hors de son périmètre : 16~s pour 1\,400 stages hors charge, dont 11,5~s dans la
feuille « Stages » pour 20~requêtes SQL, et 41 à 57~s pendant la montée en charge. Le temps se
passait dans la bibliothèque Excel de Moodle (\code{MoodleExcelWorkbook}, sur PhpSpreadsheet), qui
construit tout le classeur en mémoire et applique un style à chaque cellule écrite.

La page \code{export.php} écrit désormais le classeur en flux avec OpenSpout, livré par Moodle
depuis la version~4.3, comme la branche Excel de l'export pré-stage
(\code{request_answers_export.php}). Les trois feuilles, leurs colonnes, l'ordre des lignes et les
valeurs sont conservés.

\section{Ce qui change dans la page}

\begin{itemize}
  \item \textbf{Écriture.} Chaque ligne est envoyée au navigateur dès qu'elle est prête
    (\code{\OpenSpout\Writer\XLSX\Writer}, une feuille par
    \code{addNewSheetAndMakeItCurrent()}), avec un dossier temporaire propre à la requête
    (\code{make_request_directory()}). La session est libérée avant l'envoi
    (\code{\core\session\manager::write_close()}) : un long export ne bloque plus les autres
    onglets de l'utilisateur.
  \item \textbf{Cellules typées.} Nombres en \code{NumericCell}, textes en \code{StringCell},
    dates en \code{DateTimeCell}, jamais de \code{FormulaCell}. Les formats de
    \code{MoodleExcelWorksheet} sont repris : texte (\code{@}) pour les textes, général pour les
    nombres, en-têtes en gras. Le format texte est porté par la ligne plutôt que par chaque
    cellule : OpenSpout fusionne à chaque cellule son style avec celui de sa ligne, ce qui
    triplait le temps d'écriture des textes.
  \item \textbf{Dates.} La date et l'heure sont celles du fuseau de l'utilisateur, comme avec
    \code{write_date()}. Une étape qui n'a pas eu lieu reste une cellule de texte vide.
  \item \textbf{Protection des formules.} Sous Moodle~4.5, \code{write_string()} préfixe d'une
    apostrophe les textes commençant par \code{=}, \code{+}, \code{-} ou \code{@}
    (\code{\core\dataformat::escape_spreadsheet_formula()}). La page fait de même lorsque cette
    méthode existe, et ne le fait pas sous Moodle~4.3, qui ne la connaît pas. Les valeurs restent
    donc celles de l'ancien export sur chaque version. La restauration
    (\code{global_export_importer}) retire ce préfixe.
  \item \textbf{Réponses lues en flux.} La feuille « Réponses aux questionnaires » lit les
    réponses par un \textit{recordset}, stage après stage, avec la même requête et le même tri.
    Les notes des évaluations textuelles importées de StageVet sont calculées d'avance et
    intercalées par numéro de stage, comme le faisait le tri final de l'ancienne version. Le nom
    de l'étudiant, la thématique et l'intitulé de chaque question ne sont mis en forme qu'une fois.
\end{itemize}

\section{Méthode}

Moodle~4.5 (\code{MOODLE_405_STABLE}), PHP~8.3, PostgreSQL~16, serveur web intégré de PHP, un
seul client. L'ancienne page est copiée sous \code{export_old.php} dans le Moodle de test, pour
comparer les deux versions sur les mêmes données. Les scripts sont dans
\code{docs/audit/questions/}, les résultats bruts dans \code{logs/export.txt}.

\begin{itemize}
  \item \code{setup.php 1500}, puis \code{export_data.php} : 3\,000 stages. Chaque colonne des
    trois feuilles porte une valeur : réponses des huit types de question, dont une sur dix hors
    de la forme attendue ; détail de convention complet ; étapes datées ; deux plages pour un
    stage sur cinq ; rapports déposés ; validations et annulations ; évaluations StageVet notées,
    dont certaines sur des stages sans aucune réponse ; textes commençant par \code{=}, \code{+},
    \code{-} ou \code{@} ; dates de naissance antérieures à 1970.
  \item \code{volume.php} ajoute ensuite 5\,000 demandes de stage avec leurs réponses : 8\,000
    stages en tout.
  \item \code{export_bench.py} télécharge l'export par HTTP, dans la session d'un compte DEVE. Le
    serveur consigne la durée et le pic de mémoire PHP de chaque requête
    (\code{export_prepend.php}).
  \item \code{export_compare.py} compare les deux classeurs avec openpyxl, cellule par cellule :
    noms et ordre des feuilles, valeur et type de chaque cellule (texte, nombre, date), puis format
    de nombre et gras. Une cellule vide et un texte vide sont tenus pour égaux, de même qu'un
    nombre écrit \code{1} ou \code{1.0} : le format xlsx n'a qu'un type numérique.
\end{itemize}

\section{Résultats}

\begin{center}
\begin{tabular}{@{}lrrrr@{}}
  \toprule
  & \multicolumn{2}{c}{Durée} & \multicolumn{2}{c}{Pic de mémoire} \\
  \cmidrule(lr){2-3}\cmidrule(l){4-5}
  Stages (réponses) & avant & après & avant & après \\
  \midrule
  3\,000 (46\,200) & 77,8~s & 7,6~s & 554~Mo & 69~Mo \\
  8\,000 (86\,200) & 166,4~s & 15,4~s & 1\,089~Mo & 116~Mo \\
  \bottomrule
\end{tabular}
\end{center}

Durées : meilleure de trois exécutions, sauf l'ancienne version à 8\,000 stages (une seule).
L'export est dix fois plus rapide et occupe huit à neuf fois moins de mémoire. La nouvelle version
fait 58~requêtes en tout.

\paragraph{Équivalence.} Les classeurs ont les mêmes feuilles et les mêmes valeurs : 767\,599
cellules comparées à 3\,000 stages, 1\,602\,599 à 8\,000, sans écart. Le même résultat est obtenu en
substituant au banc l'OpenSpout~4.15 de Moodle~4.3. La restauration
(\code{global_export_importer::read()}) relit les deux classeurs à l'identique : 3\,000
enregistrements, aucun avertissement.

\paragraph{Où passe le temps.} À 3\,000 stages, sur 7,6~s : 0,3~s pour le bilan de promotion,
1,0~s pour la feuille « Stages », 0,5~s pour la requête des réponses et les notes textuelles,
3,1~s pour la feuille « Réponses » (47\,000 lignes) et 3,0~s pour la compression du fichier
(44~Mo de XML). Il reste donc de l'ordre de 60~µs par ligne de réponse dans OpenSpout lui-même.

\section{Écarts de présentation}

Deux différences de présentation subsistent ; aucune ne touche une valeur.
\begin{itemize}
  \item \textbf{Format des dates.} L'ancien export demandait le format \code{dd/mm/yyyy}, mais
    \code{MoodleExcelFormat} l'ignorait en silence (il n'accepte qu'une liste de formats
    prédéfinis) : les dates s'affichaient avec l'heure, au format \code{m/d/yy h:mm}. Elles
    s'affichent désormais en jj/mm/aaaa ; l'heure reste dans la valeur.
  \item \textbf{Retour à la ligne.} OpenSpout active le renvoi à la ligne automatique dans les
    cellules dont le texte contient un saut de ligne (adresses, commentaires, évaluations).
\end{itemize}

\section{Scénarios, montée en charge et tests unitaires}

\paragraph{Scénarios.} Les 68 scénarios de \code{scenarios.py} réussissent, dont les six contrôles
de l'export général (colonnes, réponses présentes, nombres, oui/non, cases, notes). Résultat :
\code{logs/scenarios-flux.txt}.

\paragraph{Montée en charge.} Les trois vagues de l'audit sont rejouées dans le même ordre :
site neuf de 700 étudiants (\code{setup.php 700}, 1\,400 stages), puis 16, 32 et 64 clients sur
200 stages chacune, serveur à 16 processus, deux comptes DEVE qui téléchargent en boucle l'export
général et l'export pré-stage. Résultats : \code{logs/load-flux-*.json}, à comparer à
\code{logs/load-*.json} de l'audit.

\begin{center}
\begin{tabular}{@{}rrrrrr@{}}
  \toprule
  & \multicolumn{2}{c}{Export général, médiane} & \multicolumn{2}{c}{Export général, max.} & \\
  \cmidrule(lr){2-3}\cmidrule(lr){4-5}
  Clients & audit & en flux & audit & en flux & Erreurs \\
  \midrule
  16 & 45,1~s & 6,8~s & 45,1~s & 10,7~s & 0 \\
  32 & 41,4~s & 13,2~s & 52,0~s & 14,0~s & 0 \\
  64 & 47,5~s & 14,2~s & 57,3~s & 18,8~s & 0 \\
  \bottomrule
\end{tabular}
\end{center}

Pendant les vagues, les 11 exports généraux ont pris de 2,8 à 18~s côté serveur, avec un pic de
mémoire de 36~Mo. Aucune erreur sur les 3\,809 actions des trois vagues ; l'état final est celui
de l'audit (200 stages évalués par vague, 6\,400 réponses) et \code{check.php} ne trouve aucune
réponse invalide ni aucun doublon parmi les 19\,200 réponses. Les autres actions ont des médianes
de 10 à 34~\% plus élevées que dans l'audit, sur une autre machine ; les vagues durent moins
longtemps à 32 et 64 clients (33,5 et 36,8~s contre 54,2 et 60,4~s), plus à 16 (34,7~s contre 29,4~s). Quelques maxima isolés atteignent 14 à 20~s (évaluation
de l'enseignant ou du maître de stage à 32 et 64 clients), quand un export occupe l'un des 16
processus du serveur ; l'audit en avait relevé de 52 à 58~s.

\paragraph{Ordre des opérations.} Une première série, jouée sur le site des scénarios, a vu
l'action « étudiant » refusée pour 91 à 115 stages sur 200. Les scénarios réenregistrent une liste
d'évaluation par \code{evallist_edit.php}, dont le formulaire lisait les options en
\code{PARAM_TEXT} : l'option \code{<b>Urgences</b>} y devenait \code{Urgences}, alors que le client
de charge envoie le texte d'origine. Le défaut, antérieur à l'export, touchait aussi les vrais
utilisateurs : réenregistrer une liste retirait les balises de ses options, et les réponses déjà
enregistrées avec la balise ne correspondaient plus à aucune option. Il est corrigé
(section~\ref{sec:corrections}) ; scénarios puis vague de 64~clients sur le même site ne donnent
plus aucune erreur.

\paragraph{Tests unitaires.} PHPUnit, sur le même Moodle~4.5 : \code{mod_stagesynthesis}
réussit (8~tests). \code{mod_stage} comptait 275~tests, dont 5 échecs et 1 erreur, dans des tests
écrits avant l'audit et jamais exécutés jusque-là ; aucun n'inclut \code{export.php}. Après les
corrections ci-dessous, les 276~tests réussissent (1\,653~assertions).

\section{Corrections}\label{sec:corrections}

\begin{itemize}
  \item \textbf{Options des listes d'évaluation.} Le formulaire des listes les lit désormais en
    \code{PARAM_RAW}, comme les réponses des questions à choix : elles sont enregistrées telles
    que saisies. Le seul affichage qui ne les échappait pas, la liste déroulante du questionnaire
    de demande (le gabarit Moodle y écrit le libellé sans échappement), les passe maintenant par
    \code{s()}. Un test vérifie le type des champs et l'échappement.
  \item \textbf{Réponses texte de la demande de stage.} Les espaces aux extrémités sont de nouveau
    retirés, comme avant l'arrivée des types de question ; les commentaires des évaluations
    gardent les leurs. C'est le comportement que vérifiait
    \code{request_validation_test} (\code{' Urgences '} enregistré \code{Urgences}).
  \item \textbf{Services web.} \code{get_my_stages} et \code{register_entries} emploient les
    classes \code{core_external} (Moodle~4.2 et suivants) au lieu d'inclure
    \code{lib/externallib.php}, que Moodle interdit dans un test non isolé.
  \item \textbf{Tests mis à jour.} Une réponse vide à une question obligatoire est refusée comme
    manquante (\code{errorrequiredanswer}), ce qu'elle a toujours été ; les gabarits de convention
    des tests sont de vrais PDF d'une page produits par TCPDF, et non une simple entête
    \texttt{\%PDF-1.4} que le contrôle de lisibilité refuse ; la date comparée est écrite sans
    retrait du zéro du jour (\code{userdate()} avec \texttt{\$fixday} à faux).
\end{itemize}

\section{Limites}

\begin{itemize}
  \item Le classeur étant envoyé au fil de l'eau, une erreur survenant après le début de l'envoi
    donne un fichier tronqué plutôt qu'une page d'erreur. L'export pré-stage a le même
    comportement.
  \item La mémoire reste proportionnelle au nombre de stages : les saisies, leurs conventions et
    le bilan de promotion sont toujours chargés en une fois. La limite de mémoire relevée de la
    page est conservée.
  \item La compatibilité avec Moodle~4.3 a été vérifiée pour OpenSpout seulement, sur un
    Moodle~4.5 ; la page n'emploie pas de syntaxe postérieure à PHP~8.0.
  \item Le serveur de test est le serveur intégré de PHP, comme pour l'audit ; un serveur de
    production (PHP-FPM) répartit autrement les requêtes longues.
\end{itemize}

\end{document}
